% 作者简介供本人确认；本次未取得可确认授权的两位作者照片。
\par\needspace{62mm}
\section*{作者简介}
\noindent
\begin{minipage}[t]{0.20\linewidth}
\vspace{0pt}\centering\authorphoto{fan_ju}{范举}
\end{minipage}\hfill
\begin{minipage}[t]{0.76\linewidth}
\vspace{0pt}
\textbf{范举}，中国人民大学信息学院教授、博士生导师，国家级青年人才，RUC-DataLab团队负责人。
主要研究数据治理技术与系统、智能数据库系统及数据与人工智能的交叉问题，关注数据准备、数据质量管理和数据科学智能体。
主持国家自然科学基金优秀青年基金、重点项目等，致力于将语义理解、自动化处理与可靠的数据系统结合，提升数据的可用性与应用价值。
\end{minipage}
\par\vspace{17pt}\needspace{48mm}
\noindent
\begin{minipage}[t]{0.20\linewidth}
\vspace{0pt}\centering\authorphoto{fan_meihao}{范梅浩}
\end{minipage}\hfill
\begin{minipage}[t]{0.76\linewidth}
\vspace{0pt}
\textbf{范梅浩}，中国人民大学信息学院2023级直博生，师从范举教授。
主要研究面向数据分析的大语言模型与智能体技术，重点关注任务感知的数据准备、表格理解及多步骤流程的自动生成与验证。
作为第一作者完成AutoPrep和DeepPrep等工作，相关成果发表于VLDB、ICDE等国际会议。
当前关注多模态数据准备以及数据处理过程与下游人工智能任务之间的协同优化。
\end{minipage}
